\documentclass[10pt,twocolumn]{article}
\usepackage[a4paper,margin=0.55in,columnsep=0.22in]{geometry}
\usepackage{amsmath,amssymb,graphicx,booktabs,hyperref,xcolor,titlesec,enumitem}
\usepackage[font=footnotesize,labelfont=bf,skip=2pt]{caption}
\hypersetup{colorlinks=true,urlcolor=blue!60!black,linkcolor=black}
\titlespacing*{\section}{0pt}{4pt}{2pt}
\titleformat{\section}{\normalsize\bfseries}{\thesection}{0.5em}{}
\setlist{nosep,leftmargin=*}
\setlength{\parskip}{1pt}

\title{\vspace{-0.6in}\textbf{FrictionRL: Learning When to Ask for an OTP}\\[1pt]
\large Reinforcement Learning for Adaptive Step-Up Authentication in Digital Payments}
\author{Tanisha Ray (B24CM1061) \quad Cherry Aggarwal (B24CM1016)\\
\small CSL3090 Artificial Intelligence, IIT Jodhpur \quad
GitHub: \url{https://github.com/<USER>/frictionrl}}
\date{\vspace{-0.35in}}

\begin{document}
\maketitle

\section{Problem}
Every UPI or card payment receives one of three decisions: \emph{approve}, \emph{step-up} (OTP), or \emph{block}.
Static rules face a trade-off: lax rules let a fraudster who has taken over an account drain it; strict rules
bury genuine users in OTPs and false blocks until they leave the platform. The decision is \textbf{sequential}:
(i) friction compounds---each OTP raises the chance that the next one makes a genuine user quit;
(ii) stopping a fraudster early prevents all of their later transactions;
(iii) an OTP \emph{gathers information}---a failed OTP is strong evidence of takeover;
(iv) the attacker adapts---a blocked fraudster retries with a smaller amount.
A per-transaction classifier optimises none of these. We formulate a user session as an MDP and optimise
expected cumulative \emph{rupees}.

\section{RL Formulation}
\textbf{Episode.} One user session of up to $H{=}20$ transactions. With probability $0.15$ the account is taken
over at a uniformly random step $t_0\in[2,14]$; the fraud flag is \emph{hidden}. The fraudster later cashes out
(w.p.\ 0.3 per transaction) or gives up when stopped.

\textbf{State space} $\mathcal{S}\subset\mathbb{R}^{11}$: current transaction (log of amount relative to the
user's usual amount; new device; new payee; odd hour; burst), session context (fraction of risky transactions so
far; progress $t/H$) and the agent's own history (\#OTPs issued, OTP-failed seen, OTP-passed seen, \#blocks).
The history features act as sufficient statistics for the hidden variables (fraud presence, user patience),
making the observation approximately Markov---the true problem is a POMDP.
For tabular methods we aggregate to $|\mathcal{S}_d|{=}432$ states (amount$\times$flags$\times$\#OTP$\times$
failed$\times$passed$\times$risky$\times$blocked).

\textbf{Action space} $\mathcal{A}=\{\textsc{Approve},\textsc{OTP},\textsc{Block}\}$ (discrete, $|\mathcal{A}|{=}3$).

\textbf{Reward} (rupees, scaled by $10^{-3}$ for training), for transaction amount $x$:
\begin{equation*}\scriptsize
\begin{aligned}
R ={}& 0.02x\,\mathbf{1}_{\text{gen.\ done}} - (x{+}200)\,\mathbf{1}_{\text{fraud done}} - 2\,\mathbf{1}_{\text{OTP}}
- 50\,\mathbf{1}_{\text{gen.\ blocked}}\\
&- 3000\,\mathbf{1}_{\text{churn}} - 1500\,\mathbf{1}_{\text{gen.\ frozen}}.
\end{aligned}
\end{equation*}
There is deliberately \emph{no bonus for catching fraud}: only money counts, so ``block everything'' is not rewarded.
$\gamma{=}0.99$; churn and freezing (2nd block) are absorbing terminal states.

\textbf{Dynamics $T$} (simulator). Transaction content comes from a source: a synthetic generator with deliberately
overlapping genuine/fraud distributions (default) or real feature rows replayed from IEEE-CIS (customers rebuilt
with the \texttt{card1+addr1+D1} uid). Reactions---which no dataset records---are simulated:
genuine OTP pass $0.97$; churn after OTP $=0.01+0.04\cdot\#\text{prev OTPs}$; churn after a false block
$=0.2+0.2\cdot\#\text{prev blocks}$; fraudster beats OTP $0.15$ (SIM-swap); quits after a failed OTP $0.7$;
retries at half the amount after a block $0.5$.

\section{Algorithms}
\begin{itemize}
\item \textbf{Baselines}: always-approve/OTP, random, amount threshold (swept), and a propositional rule base.
\item \textbf{Model-based learning + VI/PI}: count $(s,a,s')$ to estimate $\hat T,\hat R$, solve with value
iteration and policy iteration; iterate data collection with the new policy (3 rounds).
\item \textbf{Tabular Q-learning}: $Q(s,a)\!\leftarrow\!(1{-}\alpha)Q(s,a)+\alpha[r+\gamma\max_{a'}Q(s',a')]$,
$\alpha_n{=}n^{-0.6}$, with $\epsilon$-greedy (decaying) or exploration function $f(u,n){=}u{+}k/n$.
\item \textbf{Approximate Q}: $Q(s,a){=}\mathbf{w}_a^\top\phi(s)$, $w_i\!\leftarrow\!w_i{+}\alpha\,\delta\,f_i(s,a)$.
\item \textbf{DQN} (replay buffer 100k, target net synced every 1k steps, $[64,64]$ MLP) and \textbf{PPO}
(2048-step rollouts, GAE $\lambda{=}0.95$, clip $\epsilon{=}0.2$), 500k steps each, Stable-Baselines3.
\end{itemize}

\small
\begin{thebibliography}{9}\setlength{\itemsep}{0pt}
\bibitem{aima} S.~Russell, P.~Norvig. \emph{Artificial Intelligence: A Modern Approach}, 4th ed., 2021.
\bibitem{dqn} V.~Mnih et al. Human-level control through deep RL. \emph{Nature}, 2015.
\bibitem{ppo} J.~Schulman et al. Proximal policy optimization algorithms. arXiv:1707.06347, 2017.
\bibitem{sb3} A.~Raffin et al. Stable-Baselines3. \emph{JMLR}, 2021.
\bibitem{ieee} IEEE-CIS Fraud Detection, Kaggle, 2019.
\end{thebibliography}
\end{document}
